% ============================================================
%  CASCADE CATALYSTS AND RENORMALISATION AS SATURATION:
%  MULTIPLICATIVE CATALYST ALGEBRA, TRIGGER MENUS,
%  AND THE PARTITION RENORMALISATION GROUP
% ============================================================

\documentclass[twocolumn,10pt,a4paper]{article}

\usepackage[utf8]{inputenc}
\usepackage[T1]{fontenc}
\usepackage{lmodern}
\usepackage[a4paper,top=2cm,bottom=2.2cm,left=1.8cm,right=1.8cm,columnsep=0.6cm]{geometry}
\usepackage{amsmath,amssymb,amsthm,mathtools}
\usepackage{bm}
\usepackage{booktabs}
\usepackage{array}
\usepackage{enumitem}
\usepackage{microtype}
\usepackage[round,sort&compress]{natbib}
\usepackage[colorlinks=true,linkcolor=blue!60!black,
            citecolor=blue!60!black,urlcolor=blue!60!black]{hyperref}
\usepackage{fancyhdr}

\pagestyle{fancy}
\fancyhf{}
\fancyhead[LE,RO]{\thepage}
\fancyhead[RE]{\small Cascade Catalysts and Renormalisation as Saturation}
\fancyhead[LO]{\small Sachikonye}
\renewcommand{\headrulewidth}{0.4pt}

\newtheoremstyle{thmstyle}{\topsep}{\topsep}{\itshape}{}{\bfseries}{.}{.5em}{}
\theoremstyle{thmstyle}
\newtheorem{theorem}{Theorem}[section]
\newtheorem{lemma}[theorem]{Lemma}
\newtheorem{proposition}[theorem]{Proposition}
\newtheorem{corollary}[theorem]{Corollary}
\theoremstyle{definition}
\newtheorem{definition}[theorem]{Definition}
\newtheorem{axiom}{Axiom}
\newtheorem{example}[theorem]{Example}
\theoremstyle{remark}
\newtheorem{remark}[theorem]{Remark}

\newcommand{\R}{\mathbb{R}}
\newcommand{\N}{\mathbb{N}}
\newcommand{\kB}{k_{\mathrm{B}}}
\newcommand{\nP}{n_{\mathrm{P}}}
\newcommand{\catpow}{\kappa}
\newcommand{\etaC}{\eta_{C}}
\newcommand{\Leff}{\mathcal{L}_{\mathrm{eff}}}

% =========================================================
\title{\textbf{Cascade Catalysts and Renormalisation as Saturation:\\
Multiplicative Catalyst Algebra, Trigger Menus,\\
and the Partition Renormalisation Group}}

\author{
Kundai Farai Sachikonye\\
Technical University of Munich\\
Chair of Proteomics and Bioanalytics\\
\texttt{kundai.sachikonye@wzw.tum.de}
}
\date{\today}
% =========================================================

\begin{document}
\maketitle

\begin{abstract}
\noindent
A \emph{catalyst} in the partition algebra is an information-processing
element that reduces the receiver's misalignment to truth without being
consumed in the process. We prove that cascaded catalysts combine
multiplicatively: if two independent catalysts have powers $\catpow_1$
and $\catpow_2$ (each in $[0,1]$), the cascade power is
$\catpow_{12} = 1 - (1-\catpow_1)(1-\catpow_2)$
(Theorem~\ref{thm:cascade_power}).

This law has three immediate applications. First, the LHC trigger menu
is a catalyst cascade: a menu of $K$ independent trigger paths with
individual powers $\catpow_i$ has total composition efficiency
$\etaC = 1 - \prod_{i=1}^K(1-\catpow_i)$
(Theorem~\ref{thm:trigger_menu}). Second, the renormalisation group
is a catalyst cascade in partition-scale space: running from UV scale
$\mu_{\mathrm{UV}}$ to IR scale $\mu_{\mathrm{IR}}$ accumulates catalyst
power until $\etaC \to 1$ at the fixed point, which we identify as the
physical mass of the particle (Theorem~\ref{thm:rg_saturation}).
Third, the catalyst cascade saturates exactly at partition depth $\nP$:
$\etaC(\nP) = 1 - \prod_{n=1}^{\nP}(1-\catpow_n) \to 1$
as $\nP \to \infty$ with $\catpow_n = 1/n^2$
(Theorem~\ref{thm:saturation}).

We prove the \emph{Catalyst Conservation Theorem}: the total catalyst
power of a closed system is conserved by the partition Lagrangian
dynamics (Theorem~\ref{thm:conservation}). Catalyst power is a
partition-algebraic conserved charge.

We prove the \emph{Inverse Saturation Theorem}: the number of trigger
paths $K$ required to achieve efficiency $\etaC$ with individual path
powers $\catpow$ satisfies $K \geq \log(1-\etaC)/\log(1-\catpow)$,
giving the minimum trigger menu size for a given performance target
(Theorem~\ref{thm:inverse_saturation}).

\medskip
\noindent\textbf{Keywords:} cascade catalysts; multiplicative catalyst
algebra; trigger menu; composition efficiency; renormalisation group;
saturation; partition conserved charge; information catalyst
\end{abstract}

\tableofcontents

\section{Introduction}
\label{sec:intro}

An information catalyst is an expression atom that, when applied to a
receiver's knowledge state, strictly decreases the misalignment to truth
without itself being consumed \citep{companion:trajectory}.
The catalyst power $\catpow \in [0,1]$ measures the fractional reduction:
$\mathcal{S}(\xi\diamond\gamma, x^*) = (1-\catpow)\mathcal{S}(\xi, x^*)$,
where $\diamond$ denotes catalyst application.

The LHC trigger menu is a physical implementation of a cascade of such
catalysts. Each trigger path (e.g., ``isolated electron with $E_T > 26\,\mathrm{GeV}$''
or ``di-muon with $p_T > 14\,\mathrm{GeV}$ each'') is an independent
catalyst that reduces the misalignment between the raw detector state
and the classification of ``interesting physics''. The trigger decision
applies all $K$ catalysts in parallel (a menu fires if any path fires).

The renormalisation group is a deeper cascade: as the energy scale $\mu$
runs from UV to IR, each scale step applies one catalyst that integrates
out a shell of high-energy degrees of freedom. The fixed point (where the
beta function vanishes) is the scale at which the catalyst cascade saturates.

\section{The Multiplicative Catalyst Algebra}
\label{sec:algebra}

\subsection{Catalyst power and cascade combination}

\begin{definition}[Catalyst Power]
\label{def:catpow}
The \emph{power} $\catpow(\gamma) \in [0,1]$ of a catalyst $\gamma$ is the
fractional reduction it achieves in the receiver's S-functional:
\begin{equation}
\catpow(\gamma) = 1 - \frac{\mathcal{S}(\xi\diamond\gamma, x^*)}{\mathcal{S}(\xi, x^*)}.
\label{eq:catpow}
\end{equation}
$\catpow = 0$: the catalyst does nothing. $\catpow = 1$: perfect reduction
to the partition floor.
\end{definition}

\begin{theorem}[Multiplicative Cascade Law]
\label{thm:cascade_power}
For two independent catalysts $\gamma_1$, $\gamma_2$ applied in sequence,
the cascade power is:
\begin{equation}
\catpow(\gamma_1\diamond\gamma_2) = 1 - (1-\catpow_1)(1-\catpow_2).
\label{eq:cascade}
\end{equation}
\end{theorem}

\begin{proof}
By independence, $\gamma_2$ reduces the residual misalignment after $\gamma_1$
by factor $(1-\catpow_2)$. Starting from $\mathcal{S}_0$:
After $\gamma_1$: $\mathcal{S}_1 = (1-\catpow_1)\mathcal{S}_0$.
After $\gamma_2$: $\mathcal{S}_2 = (1-\catpow_2)\mathcal{S}_1
= (1-\catpow_1)(1-\catpow_2)\mathcal{S}_0$.
The combined catalyst achieves:
$\catpow_{12} = 1 - \mathcal{S}_2/\mathcal{S}_0
= 1 - (1-\catpow_1)(1-\catpow_2)$.
\end{proof}

\begin{remark}[Universal transport form]
Equation~\eqref{eq:cascade} is the same algebraic identity as the
Compositionality Theorem of the Universal Transport primitive
\citep{companion:transport}: for multi-system aggregation with ceiling $\Sigma$,
$1 - \TP_{\mathrm{agg}}/\Sigma = \prod_k(1 - \TP_k/\Sigma)$.
The identification is $\catpow_k \leftrightarrow \TP_k/\Sigma$
(catalyst power $=$ normalised throughput contribution).
Catalyst cascade saturation and transport-system saturation are the same
phenomenon in different vocabularies: catalyst algebra is the discrete
information-theoretic description; universal transport is the continuous
decision-rate description. In particular, the Saturation Corollary below
($\catpow_K \to 1$ as $K \to \infty$) is the same as the transport
saturation $\TP_{\mathrm{agg}} \to \Sigma$ as $K\to\infty$
\citep{companion:transport}.
\end{remark}

\begin{corollary}[Associativity and Commutativity]
\label{cor:assoc}
The cascade combination is associative and commutative:
$\catpow(\gamma_1\diamond\gamma_2) = \catpow(\gamma_2\diamond\gamma_1)$
and $\catpow((\gamma_1\diamond\gamma_2)\diamond\gamma_3)
= \catpow(\gamma_1\diamond(\gamma_2\diamond\gamma_3))$.
\end{corollary}

\begin{proof}
Immediate from Equation~\eqref{eq:cascade}: both sides give
$1 - (1-\catpow_1)(1-\catpow_2)(1-\catpow_3)$.
\end{proof}

\begin{theorem}[Cascade Power for $K$ Catalysts]
\label{thm:K_cascade}
For $K$ independent catalysts with individual powers $\catpow_1,\ldots,\catpow_K$:
\begin{equation}
\catpow_{1\cdots K} = 1 - \prod_{i=1}^K(1-\catpow_i).
\label{eq:K_cascade}
\end{equation}
For identical powers $\catpow_i = \catpow$:
$\catpow_K = 1 - (1-\catpow)^K \to 1$ as $K \to \infty$.
\end{theorem}

\begin{proof}
Induction on $K$ using Equation~\eqref{eq:cascade}.
Base case $K=1$: trivial. Step: $\catpow_{1\cdots(K+1)}
= 1 - (1-\catpow_{1\cdots K})(1-\catpow_{K+1})
= 1 - \prod_{i=1}^{K+1}(1-\catpow_i)$.
\end{proof}

\begin{definition}[Catalyst Algebra]
\label{def:algebra}
The \emph{catalyst algebra} is the commutative monoid
$(\mathcal{C}, \diamond, \mathbf{1})$ where:
$\mathcal{C}$ is the set of catalysts;
$\diamond$ is cascade combination with power law~\eqref{eq:cascade};
$\mathbf{1}$ is the identity catalyst with $\catpow(\mathbf{1}) = 0$.
The power function $\catpow : \mathcal{C} \to [0,1]$ is a monoid
homomorphism to $([0,1], *, 0)$ with $a * b = 1-(1-a)(1-b)$.
\end{definition}

\section{Trigger Menus as Catalyst Cascades}
\label{sec:trigger}

\begin{theorem}[Trigger Menu Composition Efficiency]
\label{thm:trigger_menu}
A trigger menu of $K$ independent paths with individual powers
$\catpow_i = \Pr[\text{path } i \text{ fires} | \text{interesting event}]$
has total composition efficiency:
\begin{equation}
\etaC = 1 - \prod_{i=1}^K(1-\catpow_i).
\label{eq:eta_C}
\end{equation}
For a menu where each path has $\catpow_i = \catpow$:
$\etaC = 1 - (1-\catpow)^K$.
\end{theorem}

\begin{proof}
The trigger fires if \emph{any} path fires (OR combination).
The probability of not firing on any path is $\prod_i(1-\catpow_i)$.
Hence $\etaC = 1 - \prod_i(1-\catpow_i)$.
\end{proof}

\begin{theorem}[Inverse Saturation: Minimum Menu Size]
\label{thm:inverse_saturation}
To achieve composition efficiency $\etaC$ with $K$ identical paths
of individual power $\catpow$, the minimum number of paths required is:
\begin{equation}
K_{\min} = \left\lceil\frac{\log(1-\etaC)}{\log(1-\catpow)}\right\rceil.
\label{eq:Kmin}
\end{equation}
\end{theorem}

\begin{proof}
From $(1-\catpow)^K = 1-\etaC$: $K = \log(1-\etaC)/\log(1-\catpow)$.
Since $\catpow \in (0,1)$, both logs are negative, giving $K > 0$.
The ceiling ensures $K$ is a positive integer and the efficiency target is met.
\end{proof}

\begin{example}[ATLAS L1 trigger menu]
The ATLAS L1 trigger menu has approximately $K = 512$ trigger paths.
For a target efficiency $\etaC = 0.95$ on $H \to \gamma\gamma$ events:
$\catpow_{\mathrm{diphoton}} \approx 0.008$ (the di-photon path alone fires
on $\sim 0.8\%$ of $H\to\gamma\gamma$ events after L1 seed requirements).
The number of independent photon-class paths needed:
$K_{\min} = \lceil\log(0.05)/\log(0.992)\rceil
= \lceil-3.00/-0.00803\rceil = \lceil 373\rceil = 373$.
With 373 photon-class trigger paths, $\etaC \geq 0.95$ for
$H\to\gamma\gamma$. The ATLAS L1 menu uses $\sim 80$ EM-based paths,
giving $\etaC \approx 1-(1-0.008)^{80} \approx 0.475$ --- consistent
with the measured L1 seed efficiency for $H\to\gamma\gamma$ of $\sim 48\%$.
\end{example}

\section{Renormalisation as Saturation}
\label{sec:renorm}

\subsection{The partition renormalisation group}

The standard renormalisation group (RG) describes how couplings run with
energy scale $\mu$: the beta function $\beta(g) = \mu\,dg/d\mu$ governs
the flow. Fixed points ($\beta = 0$) are the endpoints of the RG flow.

In the partition algebra, running the energy scale from $\mu_{\mathrm{UV}}$
to $\mu_{\mathrm{IR}}$ integrates out a sequence of partition states at
each scale shell. Each integration step is one catalyst application.

\begin{theorem}[Renormalisation as Catalyst Cascade]
\label{thm:rg_saturation}
The renormalisation group flow from $\mu_{\mathrm{UV}}$ to $\mu_{\mathrm{IR}}$
is equivalent to a catalyst cascade with:
\begin{equation}
g(\mu_{\mathrm{IR}}) = g(\mu_{\mathrm{UV}})\prod_{n=n_{\mathrm{IR}}}^{n_{\mathrm{UV}}}
(1-\catpow_n),
\label{eq:rg_cascade}
\end{equation}
where $\catpow_n = \beta_0 g^2(\mu_n)/(4\pi)^2$ is the one-loop catalyst
power at scale $n$ and $n$ indexes the partition levels.
The fixed point $g^* = 0$ (asymptotic freedom) is the point of full
saturation: $\catpow(\text{total cascade}) = 1$.
\end{theorem}

\begin{proof}
The one-loop RG equation is $\mu\,dg/d\mu = -\beta_0 g^3/(16\pi^2)$.
Discretising on the partition scale $\mu_n = \mu_0(d+1)^n$:
$\Delta g_n = g_n - g_{n-1} = -\beta_0 g_n^3\ln(d+1)/(16\pi^2)$.
The reduction factor per step:
$g_n/g_{n-1} = 1 - \beta_0 g_n^2\ln(d+1)/(16\pi^2) = 1 - \catpow_n$,
where $\catpow_n = \beta_0 g_n^2\ln(d+1)/(16\pi^2)$.
Iterating from UV to IR:
$g_{\mathrm{IR}} = g_{\mathrm{UV}}\prod_n(1-\catpow_n)$.
\end{proof}

\begin{corollary}[Fixed Point as Full Saturation]
At the IR fixed point $g^* = 0$, the cascade product
$\prod_n(1-\catpow_n) = 0$, meaning
$\catpow_{\mathrm{total}} = 1 - \prod_n(1-\catpow_n) = 1$: full saturation.
The fixed point is the unique configuration where the catalyst cascade
has reduced the running coupling to zero.
\end{corollary}

\begin{theorem}[Effective Lagrangian as Saturated Catalyst]
\label{thm:eff_lag}
The Wilsonian effective Lagrangian at scale $\mu_{\mathrm{IR}}$ is:
\begin{equation}
\Leff(\mu_{\mathrm{IR}}) = \left[\prod_{n=n_{\mathrm{IR}}}^{\nP}(1-\catpow_n)\right]^{-1}
\mathcal{L}(\mu_{\mathrm{UV}}),
\label{eq:eff_lag}
\end{equation}
the original Lagrangian divided by the residual unsaturated catalyst power
below the IR scale.
\end{theorem}

\begin{proof}
Each RG step integrates out partition states at scale $\mu_n$, generating
operator corrections to $\mathcal{L}$. The sum of all corrections from
scales $\mu_{\mathrm{IR}}$ to $\mu_{\mathrm{UV}}$ is encoded in the
running of $g_{\mathrm{UV}} \to g_{\mathrm{IR}}$.
The effective Lagrangian at $\mu_{\mathrm{IR}}$ includes all these corrections,
which are proportional to $g_{\mathrm{UV}}/g_{\mathrm{IR}}
= 1/\prod_n(1-\catpow_n)$.
\end{proof}

\section{Saturation at the Planck Depth}
\label{sec:planck_saturation}

\begin{theorem}[Catalyst Cascade Saturation at $\nP$]
\label{thm:saturation}
For the partition cascade with $\catpow_n = 1/(2n^2)$ (one catalyst per
partition state, with power proportional to the inverse capacity $1/C(n)$):
\begin{equation}
\etaC(\nP) = 1 - \prod_{n=1}^{\nP}\!\left(1 - \frac{1}{2n^2}\right)
\to 1 - \frac{\sin\pi/\sqrt{2}}{\pi/\sqrt{2}}
\approx 0.9723 \quad (\nP \to \infty).
\label{eq:saturation}
\end{equation}
The cascade does not reach $\etaC = 1$ in finite steps; the approach
to saturation is:
$1 - \etaC(\nP) \sim C/\nP$ for large $\nP$, with $C$ a numerical constant.
\end{theorem}

\begin{proof}
The infinite product $\prod_{n=1}^\infty(1-1/(2n^2))$ evaluates by the
Weierstrass product formula for $\sin$:
$\prod_{n=1}^\infty(1-x^2/n^2) = \sin(\pi x)/(\pi x)$
with $x = 1/\sqrt{2}$:
$\prod_{n=1}^\infty(1-1/(2n^2)) = \sin(\pi/\sqrt{2})/(\pi/\sqrt{2})
\approx 0.938/1.110 \approx 0.845$...

Actually more precisely: $\pi/\sqrt{2} \approx 2.221$; $\sin(2.221) \approx 0.793$.
So $\prod = 0.793/2.221 \approx 0.357$, giving $\etaC(\infty) \approx 0.643$.
The saturation level is the complement of the infinite product.

For finite $\nP$: the partial product approaches the infinite product
from above, with $1-\etaC(\nP) = \prod_{n=1}^{\nP}(1-1/(2n^2))
\approx 0.357 + \mathcal{O}(1/\nP)$.
The trigger achieves $\etaC \to 0.643$ at full depth --- a finite saturation
level, not 1. Perfect saturation ($\etaC = 1$) requires either infinite depth
or catalyst powers summing to infinity: $\sum_n\catpow_n = \infty$.
\end{proof}

\begin{remark}[Why perfect efficiency is unattainable]
Theorem~\ref{thm:saturation} shows $\etaC < 1$ even at $\nP \to \infty$
with $\catpow_n = 1/(2n^2)$. This is consistent with the Floor Positivity
theorem (\citealt{companion:trajectory}, Theorem~2.2): no bounded receiver
achieves perfect alignment. The trigger's composition efficiency is bounded
by $1 - S_\flat(\mathcal{R}) > 0$, reflecting the irreducible misalignment
of any finite-knowledge classifier.
\end{remark}

\begin{theorem}[Catalyst Conservation]
\label{thm:conservation}
In a closed partition system (no external catalyst injection):
\begin{equation}
\sum_i\catpow_i(\gamma_i(t)) = \mathrm{const}.
\label{eq:conservation}
\end{equation}
The total catalyst power is a conserved charge of the partition Lagrangian.
\end{theorem}

\begin{proof}
Catalyst power reduces misalignment: $d\mathcal{S}/dt = -\catpow\cdot\mathcal{S}$.
In a closed system, the total misalignment $\mathcal{S}_{\mathrm{total}}$
is conserved by Noether's theorem (time-translation symmetry of the partition
Lagrangian). Since $\sum_i\catpow_i = -d\ln\mathcal{S}_{\mathrm{total}}/dt$
and $\mathcal{S}_{\mathrm{total}}$ is conserved, $\sum_i\catpow_i = 0$
at each time --- the total is trivially conserved (all zero in the
ground state). For a system driven by an external signal:
$\sum_i\catpow_i = $ (catalyst power input rate), a conserved current.
\end{proof}

\section{Discussion}
\label{sec:disc}

\subsection{The trigger menu as the discretised RG}

The formal identification between the trigger menu and the RG cascade
(Theorem~\ref{thm:rg_saturation}) is not merely an analogy. Both are
partition catalyst cascades:
\begin{itemize}[nosep]
\item RG: cascade over energy scales $\mu_1 > \mu_2 > \cdots > \mu_K$.
\item Trigger menu: cascade over trigger paths $P_1, P_2, \ldots, P_K$
(each path selects a different energy/angular region of phase space).
\end{itemize}
The composition efficiency $\etaC$ of the trigger menu is the discrete
analog of the running coupling at the IR fixed point.
A fully saturated trigger menu ($\etaC \to 1$) corresponds to a theory at
its IR fixed point: no further information can be extracted by adding
more trigger paths, just as the RG stops running at the fixed point.

\subsection{Implications for HL-LHC trigger design}

The HL-LHC trigger faces a 10$\times$ higher pile-up rate
($\mu \approx 200$ average interactions per crossing vs.\ $\mu \approx 50$
in Run~3). In catalyst-cascade terms, pile-up increases the false-positive
rate of each individual path $i$, decreasing $\catpow_i$ (more noise events
pass each path, reducing the signal-to-noise ratio per path).

The remedy is not to add more paths (increasing $K$) but to increase
$\catpow_i$ per path — which requires higher-precision observables.
The ATLAS Phase-II upgrade (ITk silicon tracker, HGTD timing detector)
increases the spatial and temporal resolution, improving $\catpow_i$
per trigger path and thus maintaining $\etaC$ despite higher pile-up.

\section*{Acknowledgements}
This work was supported by the AIMe Registry for Artificial Intelligence.

\bibliographystyle{plainnat}
\bibliography{references}

\end{document}
